\section{Gaussian calculus on finite path trees}
\label{sec:gaussian-calculus}

We now pass from finite path calculations to the limiting Hilbert
space used in the spectral formula.  The construction must also control
the action of path matrices on vectors orthogonal to that space.
These are two different approximation problems.  Spatial square-mean
estimates determine the limiting recipes, but they cannot discard a
diagonal supported on a vanishing fraction of sites: that diagonal
may still have operator norm one.  We therefore retain such diagonal
actions in the operator representation and keep the two error norms
separate throughout the argument.

\subsection{Recipes, conditional prediction, and hidden vectors}

Begin with the spin vector, constant vectors, and a fixed countable
collection of independent $\operatorname{Uniform}[0,1]$ site seed
arrays.  A \emph{recipe} is a finite calculation from these inputs,
using multiplication by $J$, finite linear combinations, and bounded
smooth componentwise functions with bounded derivatives.  When its
output represents a direction, we include the normalization
$n^{-1/2}$.  An unbounded linear output is admitted through spatial
square-mean approximation, provided its square tails are uniformly
negligible.

The seed convention is fixed at this point.  Sample the arrays
independently of $J$ and extract them once in a typical realization
for the countably many calculations being considered.  Conditional
simulations use independent additional seeds.  Choose a countable
dense family of scalar functions and coefficients; this makes the
eventual completion separable.  Deterministic Gram limits then give
an isometry on the named recipes between completions from different
typical realizations of the same seed law.  In particular, the family
of seed generators stays fixed during the construction.

Allowing finitely many observed states and site seeds on a tree of
\Cref{def:path-tree} gives a path recipe.  A \emph{path label} is a
bounded site function obtained from such a finite calculation, from
a capped elementary likelihood, or from a successive tracial
approximation of either.  Here tracial approximation means that, with
the coordinate caps fixed,
\[
 \frac1n\sum_i|d_i-\widetilde d_i|^2\longrightarrow0
\]
in the successive limits: first dimension, then approximation
accuracy.  This controls spatial averages; it does not make the
associated diagonal matrices close in operator norm.  A
\emph{path matrix} is a fixed word, or a finite linear combination
of words, in $J$ and the corresponding diagonal labels.  The class
also contains bounded finite-rank dyads of path recipe vectors and
their successive operator-norm limits.  We call these dyads
\emph{visible}.

Every elementary likelihood has a fixed cap when it enters this
algebra.  Removing that cap later uses \Cref{prop:path-rough}; it
requires no operator-norm approximation of an unbounded likelihood.

The theorem below has two kinds of conclusion.  Its empirical and
prediction assertions construct the limiting recipe space.  Its
operator assertions control matrix words on vectors orthogonal to
that space, while keeping diagonal actions at exceptional sites.
The contraction rule in those assertions averages a returning pair
as $JDJ\mapsto\beta^2\tau(D)I$, where $\tau(D)$ is the limiting
normalized trace.  Complete noncrossing contractions repeat this
operation until no matrix edge remains; a diagonal entry retains its
outer site label instead of averaging that label.  The creation
representation and the general contraction rule are constructed below.

\begin{theorem}[Finite-tree Gaussian calculus]
\label{thm:path-calculus}
Fix $0<\beta<1$, a finite tree, a finite collection of words and
recipes, and fixed caps on its diagonal labels.  Dimension is taken
to infinity before any recipe accuracy, mesh, cap, or time limit.
The following assertions hold in the typicality convention of
\Cref{def:path-tree}.
\begin{enumerate}[label=(\roman*)]
\item Stationary recipe overlaps, normalized traces, and joint empirical
averages have deterministic limits.  The completion of one-state
recipes for the limiting inner product is a separable real Hilbert
space $\mathcal H$.  Conditional means of path recipes at a chosen
endpoint have approximations in $\mathcal H$.  The same stationary
predictors work on actual conditional trees: centered copies sharing
that endpoint have vanishing spatial pair overlaps.
\item If $u_n$ is measurable at that endpoint,
$\E\norm{u_n}^2\le C$, $\norm{u_n}\le n^M$, and
\begin{equation}
 \E|u_n^tb_n|^2\longrightarrow0
 \quad\text{for every endpoint recipe }b_n,
 \label{eq:path-hidden}
\end{equation}
then $u_n$ is orthogonal in this same squared sense to all bounded
path recipe vectors.  Multiplication by bounded path words preserves
this assertion, with the transpose placed on the other vector as
necessary.  Such vectors will be called \emph{hidden}.
\item On the complement of visible dyads, joint limiting word inner
products of hidden vectors are represented on a Hilbert space with
$J/\beta=s=\ell+\ell^*$ and
\begin{equation}
 \ell^*D\ell=\tau(D)I.
 \label{eq:path-creation}
\end{equation}
Here $\tau$ is the limiting spatial state on the diagonal algebra.
The representation retains an arbitrary bounded base action of every
diagonal, including types of zero limiting proportion.  A single
representation suffices for all the finitely many vectors and words
on the fixed tree.  No compatible choice between different trees is
asserted.
\item Normalized traces are given by complete noncrossing contractions.
For a word $Q$, let $\Delta(Q)$ be its complete contraction with the
outer diagonal unaveraged.  Then
\begin{equation}
 \max_i|Q_{ii}-\Delta(Q)_i|\longrightarrow0,
 \qquad \max_{i\ne j}|Q_{ij}|\longrightarrow0.
 \label{eq:path-entry-prediction}
\end{equation}
For two bounded path matrices,
\begin{equation}
 Q\circ R=D_{\Delta(Q)\Delta(R)}+V+o_{\op}(1),
 \label{eq:path-hadamard}
\end{equation}
where $V$ is in the operator-norm closure of bounded visible dyads.
\item Bounded clipped likelihood labels admit tracial $L^2$
approximations by fixed finite smooth multitime calculations satisfying
the mixed-slot bounds of \Cref{lem:path-mixed-tensors}, with observation
times inside their path span.  One-sided endpoint
approximations are allowed.  When only spin observations are used,
the likelihood is the flip-rate likelihood.  Predictors of these
labels mean predictors of these tracial approximants in successive
limits, with the same cap.
\item Smooth site-recipe difference matrices, their bounded products,
and matrices with column $j$ evaluated at $x^j$ belong to the same
operator closure.  Their off-diagonal difference entries vanish
uniformly.  Conditional expected derivative matrices have the stronger
pointwise truncation property in \Cref{lem:path-conditional-derivative}
for those fixed smooth calculations.  Their opened derivative terms
reduce to path matrices by \Cref{prop:path-opened-derivative}.
\end{enumerate}
For two independent stationary roots in the same disorder, the limiting
spin, field, and bounded likelihood histories at a uniformly sampled
site are independent copies.  This last assertion is restricted to
these spin-generated histories.  It is not asserted for arbitrary
noncentered recipes or recipes sharing site seeds.
\end{theorem}

The proof first constructs empirical limits and endpoint predictors,
then supplies the likelihood labels.  We next prove the norm comparison
for deterministic diagonals, including rare types, and transfer it to
path labels by completing the conditional Gaussian matrix.  Entrywise
and derivative calculations finish the operator closure.  All Gaussian
matrices below have off-diagonal variance $1/n$ unless a factor
$\beta$ is displayed.

\begin{lemma}[Planting and finite adaptive regression]
\label{lem:path-regression}
A fixed finite frozen simulation and a fixed finite list of subsequent
recipe queries have deterministic empirical limits with exponentially
small deviations at every fixed positive accuracy.  Conditional on
all the queried columns, the unrevealed Gaussian matrix is an
independent Gaussian orthogonal ensemble (GOE) matrix compressed to
their orthogonal complement.  The
revealed part is a bounded sum of dyads of calculation vectors.
\end{lemma}

\begin{proof}
We first transfer exponential estimates to a Gaussian experiment in
which the initial spin is planted.  Under the annealed one-replica
weight, the initial spin $\sigma$ is uniform and the off-diagonal
disorder has the form
\[
 J=\beta G+\frac{\beta^2}{n}\sigma\sigma^t
                     +\text{a diagonal correction},
\]
where $G$ is independent GOE.  The correction from its Gaussian
diagonal has norm tending to zero.  Its mean shift on the diagonal
has norm $O(1/n)$.  The Gibbs-to-planted comparison loses only
$e^{o(n)}$: indeed
\[
 \frac{\E Z^2}{(\E Z)^2}
 \le C_\beta\E\exp\!\left(\frac{\beta^2R_n^2}{2n}\right)
 \le C'_\beta,
\]
where $R_n$ is a sum of iid signs.  Gaussian concentration of $\log Z$
and Paley--Zygmund then give
$\log Z-\log\E Z=o(n)$ in probability.  Thus an annealed exponential
failure transfers to an exponentially small Gibbs failure on a
high-probability disorder event.  This calculation explains why the planting comparison remains valid
throughout $\beta<1$.

We can now describe exactly which randomness remains after a query.
Let $U$ have orthonormal columns, put $P=I-UU^t$, and suppose
$Y=GU$ has been revealed.  Symmetry gives $U^tY=Y^tU$.  The
conditional matrix is
\begin{equation}
 G=YU^t+UY^t-U(U^tY)U^t+PG'P,
 \label{eq:path-regression-formula}
\end{equation}
with $G'$ a fresh GOE independent of the revealed sigma-field.
To verify independence, test the GOE covariance
\[
 \E[G_{ij}G_{kl}]=n^{-1}
       (\delta_{ik}\delta_{jl}+\delta_{il}\delta_{jk})
\]
against a matrix supported in $P\R^n$ and against each revealed
column.  The covariance is zero.  Joint Gaussianity proves
independence, and the displayed revealed part is the unique symmetric
completion with those columns and zero $P$--$P$ block.

Conditional Gaussian regression along iterative spin-glass calculations
also underlies Bolthausen's construction of TAP solutions
\cite[Sections~2--3 of the arXiv version]{Bolthausen2014}.
Here it is used for finite causal
path simulations and their extra recipe queries; the required empirical
and operator conclusions are proved separately. The same formula applies
adaptively.  Given previous answers, normalize
the residual of the next query to a unit vector $q\in P\R^n$.
Its new perpendicular answer has law
$n^{-1/2}(P-qq^t)g$, and its component along $q$ has law
$\sqrt{2/n}\,zq$, with $g,z$ independent standard Gaussians.
After this answer the remaining compression is again independent.
Projection coefficients are empirical inner products of previous
calculation vectors.  Thus every fixed finite recursion reduces to
iid site Gaussians and seeds, finitely many empirical scalar
coefficients, and finitely many rank corrections.

It remains to control the finite recursion uniformly at a fixed
accuracy.  For bounded smooth tests, iid empirical concentration and
Gaussian quadratic tails propagate through each successive query.
If a query residual has length tending to zero, omit that query;
the spectral bound makes the normalized square error vanish.
Equivalently, fix a residual-length threshold, propagate its
deterministic error through the finite recursion, and send the
threshold to zero only after the dimension limit.  Fresh uniform
thresholds are continuity tests.  They can therefore be smoothed
and the smoothing subsequently removed without changing the limits.
This proves deterministic empirical convergence with exponentially
small failure at each fixed positive tolerance.
\end{proof}

\begin{lemma}[Conditional prediction and replica orthogonality]
\label{lem:path-prediction}
Assertions \textup{(i)} and \textup{(ii)} of
\Cref{thm:path-calculus} hold.  They remain valid after adjoining a
fixed finite list of bounded path-matrix operations.
\end{lemma}

\begin{proof}
We first construct the limiting inner product, then identify
conditional means in its completion.  Apply
\Cref{prop:path-simulation,lem:path-regression} to successively fixed
smooth simulations.  The comparison error in
\Cref{eq:path-simulation-error} makes their empirical limits Cauchy.
If a deviation had polynomially significant probability, conditioning
on it would cost only $O(\log n)$ extra entropy.  The same comparison
would then force the original limit on that event, a contradiction.
Apply this argument to the weight-normalized tilted laws to obtain
the weighted typical assertions.  The limiting Gram matrices are
positive semidefinite; their quotient by the null space, followed by
completion, defines $\mathcal H$.

For the second step, let $z$ be a path recipe and put
$m(X_0)=\E[z\mid X_0]$ on the stationary rooted tree.  Replace the
future by a fixed finite smooth simulation and average its seeds.
An average of $R$ independent seed copies has spatial $L^2$ error
at most $C/R$ from that simulated conditional mean.  At each fixed
$R$, it is a one-state recipe under the seed convention above.
To compare it with the actual $m$, expand the squared error into
$\norm m^2$, its cross term with the simulation, and the simulation's
squared norm.  Realize these three terms on two independent branches
sharing $X_0$, with the same nested approximation on both copies.
\Cref{prop:path-simulation}, applied while the other branch is
passive, controls each term and proves the conditional mean-square
approximation.
They also give
\[
 n^{-1}\sum_i(z_i^{(1)}-m_i)(z_i^{(2)}-m_i)\longrightarrow0
\]
for distinct branches; here the $z_i$ are written without the
normalizing factor.  Reversibility identifies stationary reverse
predictors, and \Cref{lem:path-tree-transfer} plus the finite-simulation
comparison gives the same pair centering on actual conditional trees.

The remaining assertion is orthogonality to a path recipe.  Let $u$
satisfy \Cref{eq:path-hidden} and use $R$ centered replicas
$z^{(r)}-m$.  Their Gram matrix tends, conditionally in the typical
weighted sense, to a bounded diagonal matrix.  Bessel's inequality
therefore gives
\[
 \sum_{r=1}^R|u^t(z^{(r)}-m)|^2
 \le \norm u^2\norm{\operatorname{Gram}(z^{(1)}-m,\ldots,z^{(R)}-m)}.
\]
All summands have the same expectation; divide by $R$, take dimension
large, and then send $R\to\infty$.  The $m$ term vanishes by recipe
approximation and \Cref{eq:path-hidden}.  For a matrix word $Q$,
$Q^tz$ is another path calculation with negligible square tails after
clipping its linear outputs.  The same argument applied to it proves
the last assertion.
\end{proof}

\subsection{Recovering likelihood labels from spin observations}

We next express a bounded likelihood label using finitely many spin
observations, so that it belongs to the recipe construction above.
The elementary likelihood in this subsection is the flip-rate
Radon--Nikodym derivative.  A marked-attempt likelihood is used only
when its path law or terminal projection is explicitly specified.

\begin{lemma}[Tracial likelihood approximation]
\label{lem:path-likelihood-approximation}
Assertion \textup{(v)} of \Cref{thm:path-calculus} holds, on each edge
in its original forward direction, and therefore for the joint tests
used in that theorem.
\end{lemma}

\begin{proof}
Let $c_j(x)=(1-x_j\tanh H_j)/2$ be the flip rate.  Its logarithmic
field derivatives are uniformly bounded: the first is $-x_j-m_j$,
and the next two are bounded functions of $m_j$.  For row $i$ the
field perturbation at $j$ is $u_{ij}=-2x_iJ_{ij}$ on its specified
interval mask.  The flip likelihood has logarithm
\[
 \sum_j\int\log\frac{c_j^{(i)}}{c_j}\dd N_j
                    -\sum_j\int(c_j^{(i)}-c_j)\dd s.
\]
Expand this logarithm in $J_{ij}$.  The linear term is a centered
score, the quadratic term is its bracket/compensator contribution,
and the remaining terms have absolute expectation at most
$C_T\sum_j|J_{ij}|^3=o(1)$, uniformly in $i$.  Higher moments are
controlled by \Cref{eq:path-likelihood-moment}.  Both the score
integrands and the coefficients of $J_{ij}^2$ are bounded smooth
site functions.  We may consequently cap the exponential at a fixed
continuity value, or use a smooth cap, with arbitrarily small spatial
square-mean loss.

Freeze the bounded integrands on short bins.  A bin with two or more
attempts at one site has probability $O(h^2)$; summing its squared
Poisson counts shows that ignoring all such bins costs $O_T(h)$ in
normalized mean square.  On the remaining bins the spin values at the
two ends determine whether a flip occurred.  Bounded derivatives,
\Cref{eq:path-bin-square}, count truncation, and predictable rate
averaging show that freezing the field-dependent integrands also has
vanishing normalized mean-square cost.  Multiplication by $J$
preserves that cost.

The preceding freezing estimate still leaves the outside row-spin
factor $x_i(s)$.  Freeze this factor at the beginning of its bin.
The resulting error is a vector martingale, with row-$i$ bracket
bounded by
\[
 C\int\1_{\{x_i(s)\ne x_i(\lfloor s\rfloor_h)\}}\dd s.
\]
The total bracket trace is $O_T(hn)$ on the Poisson count event, while
the operator norm of its bracket density is bounded by $C\norm J^2$.
A vector jump is bounded in Euclidean norm.  Apply the scalar
bounded-jump inequality to its squared norm, stopped at a fixed
multiple of $\sqrt n$: its jumps are $O(\sqrt n)$ and its bracket
is $O_T(n)$.  For every fixed tolerance exceeding the $O(h)$ bracket
trace, failure has superpolynomially small probability.  Thus the
fraction of rows with a significant score error tends to zero with
arbitrary polynomial accuracy.  Bracket terms need only absolute
rate averaging.  Use a finer inner mesh when approximating integrands
inside a bin whose outside factor has already been frozen.

After these error estimates, the retained score can be written as a
finite calculation from observed spins.  On a bin $[s,t]$, set
$X=X_s$, $Y=X_t$, $h=t-s$, and freeze $m=m(X)$ and $v=v(X)$.
The first log-rate derivative is $-X_j-m_j$; outside the negligible
multiple-flip bins,
\[
 \Delta N_j=(1-X_jY_j)/2,
 \qquad
 (-X_j-m_j)\Delta N_j
   =(Y_j-X_j-m_j+m_jX_jY_j)/2.
\]
Since $(-X_j-m_j)c_j=-X_jv_j/2$, the vector of frozen row scores is
\begin{equation}
 D_XJ\{X-Y+m-mXY-hXv\},
 \label{eq:path-frozen-score}
\end{equation}
where products inside braces are componentwise.  In particular the
mixed query $J(mXY)$ is retained.  It need not be a finite sum of
sitewise products of single-time outputs.  Quadratic bracket terms
are finite calculations using $J\circ J$ and bounded site factors;
this matrix has bounded operator norm and bounded absolute row sums.
Their values are therefore uniformly bounded at each fixed mesh.

For the cap of the exponential, choose a smooth bounded function
$\psi$ of the logarithmic score, with globally bounded derivatives,
which agrees with the desired clipped exponential outside an arbitrarily
small smoothing interval.  An outside spin multiplying an unbounded
query is handled by its two sign cases, not by assuming the map
$(a,u)\mapsto\psi(au)$ has bounded derivatives for arbitrary bounded
$a$ and unbounded $u$.  Explicitly,
\begin{equation}
 \psi(x_iU_i)=\sum_{\varepsilon=\pm1}
       \frac{1+\varepsilon x_i}{2}\,\psi(\varepsilon U_i).
 \label{eq:path-score-signs}
\end{equation}
For finitely many bins, split the finite sign pattern of all the frozen
outside row spins.  Each summand is a bounded product of spin selectors
and a smooth cap applied to a fixed signed linear combination of the
mixed $J$ queries and the bounded quadratic terms.  All scalar
coefficient bounds are global, independently of the query values.
Products of already bounded factors have smooth bounded extensions.
Thus \Cref{lem:path-mixed-tensors} applies to every fixed capped
calculation, including all its mixed temporal slots.

Fix the mesh, smoothing, sign decomposition, and cap before taking
the dimension limit.  When these choices are subsequently refined,
their derivative constants may grow.  The refinement removes the
spatial approximation errors established above; it makes no claim
that derivatives of those errors converge.  Observation times can
be strictly inside their bins, and can be chosen one-sidedly at
endpoints.
For predictors first take these fixed calculations, then their exact
stationary conditional means.  Conditional expectation contracts
spatial $L^2$; replicated pair tests and
\Cref{prop:path-simulation} give the one-state recipe approximation.
When an exact predictor is differentiated, use its fixed multitime
calculation on an added stationary causal branch and apply
\Cref{lem:path-conditional-derivative}.  The opened tensors are treated
by \Cref{prop:path-opened-derivative}.  No score estimate on an
artificial reverse branch with a tilted endpoint, and no derivative
estimate for a tracial approximation error, is used.
\end{proof}

\subsection{Fixed words and their Fock norm}

The creation and annihilation representation and noncrossing contractions
use standard free-probability constructions; see
Nica and Speicher~\cite[Lectures~7 and~22]{NicaSpeicher2006}.
Strong convergence of random-matrix polynomials has a related general
form in Male~\cite[Theorem~1.6]{Male2010}.  That theorem concerns
independent Gaussian unitary ensemble (GUE) matrices and an independent
auxiliary matrix family.  The auxiliary family must already converge
almost surely in normalized trace and operator norm for every
$*$-polynomial, with a faithful limiting trace.  Here the matrices are
real symmetric, and zero-frequency diagonal types remain in the norm
comparison.  We prove the required version locally, including those
types, before transferring it to adaptive path labels.

For this construction, bin the joint values of the finitely many bounded
diagonals into $q$ types.  Write $p_a$ for their limiting proportions,
set $\mathscr D=\R^q$, and put $\tau(d)=\sum_ap_ad(a)$.
For a Hilbert space $\mathcal B$ carrying a bounded representation
$\rho$ of $\mathscr D$, set
\begin{equation}
 \mathcal F_{\tau,\mathcal B}
 =\mathcal B\oplus\bigoplus_{r\ge1}
             L^2(\mathscr D,\tau)^{\otimes r}\otimes\mathcal B.
 \label{eq:path-fock-space}
\end{equation}
A diagonal acts by $\rho$ on the base and by multiplication on the
first letter elsewhere.  Define $\ell\xi=1\otimes\xi$.
Then $\ell$ is an isometry and \Cref{eq:path-creation} holds.
If $p_a=0$, the type disappears from the letter space but may act
nontrivially on $\mathcal B$.  For upper norm comparisons take a
faithful base containing every type.  The choice of multiplicity of a
base type does not affect polynomial norms.

A Wick string of degree $k$ with coefficient tensor $F$ is defined
by linearity from the words
\[
 D_{a_0}sD_{a_1}\cdots sD_{a_k},
 \qquad F=F(a_0,\ldots,a_k),
\]
by subtracting their internal adjacent returning pairs.  Equivalently, use inclusion and
exclusion over sets of disjoint adjacent edge pairs, each contraction
replacing its middle diagonal by its trace.  In the representation
\Cref{eq:path-fock-space}, its expansion consists of the $k+1$ cuts
with creations before the cut and annihilations after it.  For example,
\[
 W(a_0,a_1,a_2)
 =D_{a_0}sD_{a_1}sD_{a_2}
             -\tau(a_1)D_{a_0a_2}.
\]
An ordinary word is a sum of Wick strings obtained by noncrossing
contractions; the endpoints of each unpaired string stay unaveraged.

\begin{lemma}[Fixed-word norm comparison]
\label{lem:path-fock-norm}
Let $G_n$ be GOE, independent of a fixed finite family of deterministic
bounded diagonals with convergent joint type proportions.  For every
fixed matrix polynomial $P$ in these diagonals and $G_n$,
\begin{equation}
 \limsup_{n\to\infty}\norm{P(G_n,D)}_{\op}
 \le\norm{P(s,D)}_{\mathcal F_{\tau,\mathcal B}}
 \quad\text{in probability}.
 \label{eq:path-fock-domination}
\end{equation}
The assertion includes zero limiting type proportions with their base
actions retained.  At every fixed tolerance $0<a\le1$, its failure
probability is at most $2\exp(-cna^2/L^2)$ for sufficiently large $n$,
where $L\ge1$ depends only on the fixed word, coefficient caps, and spectral
cutoff, and not on the smallest positive type proportion.
\end{lemma}

\begin{proof}
The norm proof has three stages: types of positive proportion, rare
coordinate types, and a concentration estimate uniform in the type
counts.  We begin with the first stage.  Every finite-dimensional Wick
contraction uses the empirical trace $\tau_n=n^{-1}\Tr$, so that
internal returning pairs cancel exactly.  During the moment calculation,
$p_a$ means the empirical type proportion and $\tau$ still means
the limiting state.  Suppose first that $p_a\ge\alpha>0$ for every
type.  If every limiting proportion is positive, this holds eventually
after decreasing $\alpha$.  For a homogeneous Wick string $W_F$ of
degree $k$, use the product-state norm
\[
 \norm F_2^2=\sum_{a_0,\ldots,a_k}
             p_{a_0}\cdots p_{a_k}|F(a_0,\ldots,a_k)|^2.
\]
For $k=0$ the string is diagonal and has norm at most
$\alpha^{-1/2}\norm F_2$, so the preliminary bound below is
immediate.  Assume henceforth in the moment calculation that $k\ge1$.
Let $r$ be even and consider
$n^{-1}\E\Tr((W_F^*W_F)^{r/2})$.  Gaussian pairings glue the $rk$
sides of a polygon in either orientation.  There are $E=rk/2$
paired edges; if $V$ is the number of free index vertices, set
$e=E-V+1\ge0$.  A connected gluing contributes $n^{-e}$ times a
type average.  The two GOE covariance terms are precisely the two
possible band orientations, including their multiplicities for a loop.

Internal immediate untwisted returns cancel.  Indeed selecting a
specified pair of adjacent edges and applying their covariance gives
exactly the trace contraction subtracted in the Wick definition.
Inclusion and exclusion over disjoint pairs cancels every diagram with
at least one such internal return.  Hence every degree-one vertex of a
surviving graph lies at a boundary of one of the $r$ strings.  There are
at most $r$ leaves, and the degree identity gives
\begin{equation}
 \sum_v(\deg(v)-2)_+\le r+2e.
 \label{eq:path-map-degree}
\end{equation}

We need a diagram count whose loss grows polynomially with the word
degree.  Suppress every chain of degree-two vertices, keeping one
vertex for a pure circle.  By \Cref{eq:path-map-degree}, the remaining
core has $s\le C(r+e)$ edges.  Delete the $e$ edges outside a
spanning tree.  Reversing local cyclic orders at vertices untwists
all tree bands, and a rooted plane tree on $s-e$ edges has at most
$4^s$ encodings.  Reinsert the deleted bands in ordered corner slots.
Their endpoints, order, and twist signs have at most
$(Cs+C)^{Ce+C}$ choices.  Subdividing the core edges and choosing
the rooted boundary start adds a factor at most
\[
 Crk\binom{rk+s}{s}.
\]
Using $s\le C(r+e)$, assign the $O(r)$ core edges to the factor
$(C(1+k))^{Cr}$ and the additional $O(e)$ edges to
$(Crk)^{Ce+C}$.  Thus the number of diagrams of excess $e$ is at most
\begin{equation}
 (C(1+k))^{Cr}(Crk)^{Ce+C}.
 \label{eq:path-map-count}
\end{equation}
Loops, parallel bands, and the pure-circle case are included by the
same encoding.  No factor exponential in $rk$ remains: tree-band
twists were removed by vertex reversal, and only the $e$ extra bands
carry independent twist choices.

The count must now be combined with a bound for each type average.
Delete $e$ nontree edges.  At most $2e$ strings cross these edges;
bound each of them by
\[
 \norm F_\infty\le\alpha^{-(k+1)/2}\norm F_2.
\]
Every remaining string is a simple tree path, since a repeated tree
edge would force an internal immediate return.  To count repeated
type variables, let $c_v$ be the number of remaining string
occurrences at vertex $v$, and let $b_v$ count their endpoints there.
Each internal occurrence uses two edge incidences and each endpoint
uses one.  As each paired edge has at most two occurrences,
$2c_v-b_v\le2\deg(v)$, using the degrees of the original graph.
Thus \Cref{eq:path-map-degree} and $\sum_v b_v\le2r$ give
\[
 \sum_v(c_v-2)_+
 \le\sum_v(\deg(v)-2)_++\tfrac12\sum_v b_v
 \le2r+2e.
\]
Replace each occurrence beyond two by a fresh independent type
variable.  For a nonnegative summand the
comparison costs at most $\alpha^{-1}$ per replacement: an equality
constraint can increase a product-state integral by at most the inverse
of its smallest atom.  Each remaining variable occurs at most twice,
never twice within one tensor.  Successive coordinate
Cauchy--Schwarz therefore bounds the product integral by one
$\norm F_2$ per surviving string.  Altogether a diagram is bounded by
\begin{equation}
 n^{-e}\alpha^{-2r-(k+3)e}\norm F_2^r.
 \label{eq:path-map-weight}
\end{equation}

We can now sum the diagram bounds.  Choose an even $r$ of order
$\log n$, keeping $k,\alpha$ fixed.  Equations
\eqref{eq:path-map-count} and \eqref{eq:path-map-weight} give
\[
 n^{-1}\E\Tr((W_F^*W_F)^{r/2})
 \le (C_\alpha(1+k))^{Cr}\norm F_2^r
       (Crk)^C\sum_{e\ge0}
       \left\{\frac{(Crk)^C\alpha^{-(k+3)}}n\right\}^{e}.
\]
The sum is eventually at most two.  Markov's inequality for the trace
power, increasing the constant in $r/\log n$ if necessary, proves
\begin{equation}
 \limsup\norm{W_F}_{\op}
       \le C_\alpha(1+k)^C\norm F_2.
 \label{eq:path-wick-preliminary}
\end{equation}
The strategy of obtaining a polynomial word-length norm bound and then
removing its loss by taking powers has an antecedent in
Haagerup~\cite[Lemmas~1.3--1.4, pp.~283--286]{Haagerup1979}.
His setting is free-group convolution; the Wick/GOE diagram bound,
rare-type treatment, and adaptive-label passage here are proved locally.
To turn this preliminary estimate into the sharp comparison, decompose
a degree-$k$ polynomial into homogeneous Wick strings.  Complete
noncrossing pairing identifies the sum of their squared tensor norms
with the vacuum square norm.  Cauchy--Schwarz and
\Cref{eq:path-wick-preliminary} therefore bound the matrix norm by
$C_\alpha(1+k)^{C+1}\norm{P(s,D)}$.
Apply that bound to $(P^*P)^m$, whose degree is $2mk$.  Since
$\|(P^*P)^m\|=\|P\|^{2m}$ in both representations, it gives
\[
 \limsup\|P(G_n,D)\|_{\op}
 \le \bigl\{C_\alpha(1+2mk)^{C+1}\bigr\}^{1/(2m)}
                  \|P(s,D)\|.
\]
Letting $m\to\infty$ removes the polynomial loss.
For each fixed polynomial and each fixed power,
there are only finitely many empirical contraction coefficients, and
they converge to the coefficients formed with $\tau$.  The difference
of the resulting matrix polynomials is $o(1)$ in operator norm on
every fixed spectral ball.  Thus empirical and limiting Wick
conventions have the same norm limit.  Dimension is sent to infinity
before the power is increased; no rate for $\tau_n\to\tau$ is required.
This proves \Cref{eq:path-fock-domination} when all types have positive
proportion.

The positive-proportion case is complete.  We next keep the coordinate
space of the rare types instead of discarding it.  Write this space
as $R_n$, with dimension $m_n=o(n)$, and split
\[
 G_n=\begin{pmatrix}B_n&V_n\\V_n^*&C_n\end{pmatrix}
       \quad\text{on }R_n^\perp\oplus R_n.
\]
The internal rare block has norm $o(1)$.  Conditional on the bulk
block and its deterministic diagonals, $V_n$ has independent
$N(0,1/n)$ entries.  For every fixed bounded bulk word $A_n$,
Gaussian quadratic-form concentration, polarization, and a net of
size $\exp(O(m_n))$ imply
\begin{equation}
 \norm{V_n^*A_nV_n-\tau(A)I_{R_n}}_{\op}\longrightarrow0.
 \label{eq:path-rare-gram}
\end{equation}
Indeed a unit-vector quadratic form has variance $O(1/n)$ and fixed
accuracy tail $e^{-cn}$, while the net entropy is $o(n)$.
The normalized bulk traces converge by ordinary Wick pairing.
\Cref{eq:path-rare-gram} holds jointly for every finite list
of $A$, including $A^*A'$.

The remaining task is to turn the rare-coordinate Gram comparisons
into an operator-norm bound.  Take a subsequence realizing a limiting
norm witness and form its Hilbert ultraproduct.  The closed linear span of
$\{A V\xi:A\text{ a bulk word},\ \xi\in R\}$ is a reducing
subspace for the bulk word algebra, and
\Cref{eq:path-rare-gram} identifies it isometrically with
$L^2(\mathscr A,\tau)\otimes R$ through
$A\Omega\otimes\xi\mapsto AV\xi$.  Here $\mathscr A$ is the bulk
Gaussian/diagonal algebra, whose norm has just been controlled.
On this subspace, $V\xi=\Omega\otimes\xi$ and
$V^*(A\Omega\otimes\xi)=\tau(A)\xi$.
The orthogonal bulk complement has no coupling to $R$ and obeys the
bulk norm bound.  On the reducing sum of $R$ and its bulk cyclic
space, every return through a bulk word is exactly its vacuum
expectation, while the diagonal on $R$ retains its actual type value.
This is the representation induced from that base in
\Cref{eq:path-fock-space}.  Its polynomial norm is bounded by the
faithful-base Fock norm.  The norm on the orthogonal sum of the two
reducing pieces is the maximum of their restricted norms, proving
the rare-type assertion.

Both the positive-proportion and rare-type comparisons are now proved.
To obtain the stated exponential rate, project GOE onto the
operator-norm ball of a fixed radius $K>3$.  The bare-GOE comparison
already puts a median of $\norm{G_n}_{\op}$ below $3$ for large
$n$.  This norm is one-Lipschitz in Hilbert--Schmidt norm, so Gaussian
concentration gives
\[
 \P(\norm{G_n}_{\op}>K)\le e^{-c_Kn},\qquad c_K>0.
\]
A word of degree $k$ with
diagonal bound $M$ satisfies, by telescoping,
\[
 \norm{P(G,D)-P(G',D)}_{\op}
 \le L\norm{G-G'}_{\HS},
\]
where $L$ depends on $k,M$, the coefficients, and the spectral cutoff,
but not on $\alpha$, the type counts, or any compression rank.
The projection is nonexpansive in Hilbert--Schmidt norm.  Gaussian
concentration for the clipped polynomial, combined with its median
comparison to the limiting target, therefore gives the original
polynomial the failure bound
\[
 2e^{-c_0na^2/L^2}+e^{-c_Kn}.
\]
For $0<a\le1$, reducing the positive constant $c$ absorbs the escape
term and the prefactors into $2e^{-cna^2/L^2}$ for sufficiently large
$n$.  The preceding convergence puts the clipped median below the
limiting norm target plus any fixed slack.  These median comparisons
are uniform over type counts: otherwise take a failing sequence,
extract convergent proportions, and use the full-support or rare-type
case just proved.  One can realize this continuity on a fixed letter
space $\R^q$ with creation vector $(\sqrt{p_a})_a$ and a faithful
base; creation changes in norm by at most
$\norm{\sqrt p-\sqrt{p'}}_2$.  This completes the proof.
\end{proof}

\begin{lemma}[Sparse changes and a separate density mesh]
\label{lem:path-sparse}
The norm comparison in \Cref{lem:path-fock-norm} is valid simultaneously
for arbitrary bounded changes to a sufficiently small fraction of the
entries of finitely many binned labels, even when the changed entries
and their supports are chosen after seeing the Gaussian matrix.  The
conclusion is comparison with the representation of the changed
labels, retaining their base actions; it does not declare the change
small in operator norm.
\end{lemma}

\begin{proof}
The parameters are chosen by comparing the entropy of all allowed
changes with a fixed Gaussian failure rate.  First fix the word
family, caps, bin resolution, spectral cutoff, and operator slack
$a>0$.  For one label, changing at most $\delta n$ entries with
at most $q$ replacement choices gives
\begin{equation}
 \log N_{n,\delta}
 \le n\{H(\delta)+\delta\log q\}+o(n),
 \qquad
 H(\delta)=-\delta\log\delta-(1-\delta)\log(1-\delta).
 \label{eq:path-sparse-entropy}
\end{equation}
For finitely many labels add these entropies.  At fixed positive
$\delta$ this is a small positive multiple of $n$, \emph{not} $o(n)$.
Let $c a^2/L^2$ be the Gaussian rate from
\Cref{lem:path-fock-norm}.  Choose $\delta$ so that the sum in
\Cref{eq:path-sparse-entropy}, divided by $n$, is below one quarter
of that rate.  The union bound then gives a common exponential event
for every possible modification.  Polynomially many type-count cases
cost only $o(n)$.

This common event is enough for tracial label approximations.  Indeed,
if $n^{-1}\sum_i|d_i-\widetilde d_i|^2\le\eta^2$, then at most
$\delta n$ sites satisfy $|d_i-\widetilde d_i|>\eta/\sqrt\delta$.
Off those sites, binning and polynomial telescoping absorb the small
uniform error.  On those sites, the comparison for every modification
has already been secured by the union bound.

Only after fixing these label approximations choose a \emph{different}
mesh for the path density.  By \Cref{lem:path-two-mesh}, its specific
log-density price $\kappa$ can be made smaller than another quarter
of the same Gaussian rate.  A density tilt $e^{\kappa n+o(n)}$ leaves
a strictly positive exponential failure rate.  Then take $n\to\infty$.
The minimum positive type proportion never enters the rate: it only
enters the preliminary moment constant used to establish the median.
Finally let the sparse fraction, density cost, and operator slack tend
to zero in this nested order.
\end{proof}

\subsection{Removing the queried compression}

\begin{lemma}[Conditional Gaussian comparison]
\label{lem:path-compressed}
For fixed tracial label approximations and fixed path words, the
regression in \Cref{lem:path-regression}, followed by an independently
chosen finer density mesh, gives the norm and joint-vector comparison
in Assertion \textup{(iii)} of \Cref{thm:path-calculus}.
\end{lemma}

\begin{proof}
We must pass from the independent Gaussian comparison to the matrix
left by the path queries.  Fix the finite label calculations first.
After conditioning on the frozen path data, the spin inputs at each
mesh time are fixed, and the likelihood sees $J$ only through its
images on those input columns.  Retain an orthonormal basis of their
nonvanishing residual directions, together with the input seed columns
and the finitely many adaptive queries required by the fixed labels.
Formula \eqref{eq:path-regression-formula} leaves an independent
compressed Gaussian matrix.  The removed part is a bounded sum of
dyads of input and innovation vectors, of rank at most twice the
number of retained directions.

Small query residuals are handled before any matrix inverse is needed.
At the fixed density mesh, process the queries in causal order and
omit residuals shorter than a fixed threshold $\theta$.  Use the
orthogonal projection of each omitted query when evaluating later
fields.  This still defines a probability scheme: every refresh uses
a normalized Bernoulli probability depending on its past.  In
particular, no lower bound on the smallest eigenvalue of an
arbitrarily long mesh Gram matrix is assumed.
The omitted spin input has normalized length at most $\theta$.
The spectral bound, finite-bin count truncation and
\Cref{lem:path-two-mesh} bound the resulting log likelihood error by
$\kappa_\theta n$, with $\kappa_\theta\to0$ at the fixed mesh.
Subsequent label queries with small residual can be omitted using
normalized square-mean continuity.  Thus nearly singular regressions
cause a small density error, not an inverse-Gram operator factor.

Before applying a full-GOE word comparison, augment the residual
$PG'P$ by its missing Gaussian blocks, independently of the actual data.
Conditional on the revealed columns and fixed labels, this gives a full
independent GOE $G'$ on an enlarged probability space. The path-density
change depends only on the compressed residual; the independent added
blocks have density cost one. A full-GOE operator comparison is applied
only to this completed matrix. In particular, no unrestricted polynomial
norm domination for $PG'P$ by the full-GOE Fock norm is asserted.

We now apply the comparison on that enlarged space.  The GOE map from
standard Gaussian coordinates has Hilbert--Schmidt Lipschitz scale
$\sqrt{2/n}$.  For a Gaussian good set of probability at least $1/2$,
isoperimetry bounds the probability of Frobenius distance greater than
$r$ from the set by $2e^{-cnr^2}$.  Use as the good set the full-GOE
event for the required fixed-word comparisons in
\Cref{lem:path-fock-norm}.
Under a pointwise density price $e^{\kappa n+o(n)}$, taking
$r>C\sqrt\kappa+a$ leaves an exponential error. Telescoping a bounded
polynomial changes its operator norm by at most $Lr$. Both $L$ and the
Gaussian scale are independent of the number of queried directions;
compressing a perturbation has norm at most its original norm. The
countably many fixed approximations are taken in successive limits,
not intersected as a dimension-dependent word family. Sparse label
changes are covered by \Cref{lem:path-sparse} on this same event.

Only the hidden word inner products remain to be compared with the
actual compressed matrix.  The discrepancy between the actual and
completed matrices has range and corange in the span of the queried
directions and finitely many added isotropic Gaussian columns.
Expand each polynomial discrepancy as a sum of dyads.  Its two
vectors lie in the cyclic spaces of those directions under the actual
bounded words and labels. The queried
directions and their images are path vectors, so
\Cref{lem:path-prediction} makes them orthogonal to hidden endpoint
vectors. For an added Gaussian column $g$ of covariance at most $I/n$,
condition on all actual data. For any actual bounded word $Q$ and
actual endpoint vector $u$,
\[
 \E_g|u^tQg|^2\le n^{-1}\norm{Q^tu}^2
                    \le Cn^{-1}\norm u^2.
\]
The same estimate treats the finitely many added columns jointly and,
by a further finite polynomial expansion, their word orbits. Thus the
added cyclic directions vanish in hidden word inner products as well.
The hidden vectors do not anticipate the added blocks: these blocks
were sampled only after fixing the actual data. Removing them at this
stage proves the comparison on the hidden complement that is required
here, while retaining the rare diagonal base of the full-GOE model.

This gives a positive matrix-valued functional for every finite
family of hidden vectors on the algebra generated by $D$ and $s$,
bounded by the norm in \Cref{eq:path-fock-space}.  To represent
creation as well as the self-adjoint sum, lift the word inner products
to the algebra containing $\ell$.  Use a norm-preserving
Hahn--Banach extension of the positive functional from the unital
algebra generated by $D,s$ to the one generated by $D,\ell$.
Positivity follows from preservation of the norm and its value at the
identity.  Apply this to matrix amplifications to retain all the joint
vector states, and use the Gelfand--Naimark--Segal (GNS) construction.
The resulting lifted
vectors have exactly the original joint word inner products, while
$\ell^*D\ell=\tau(D)I$ holds in the ambient representation.
If a finite star has a permutation symmetry one may average the
extension over that finite group and take direct sums.  This proves
precisely the one-tree assertion; it imposes no symmetry joining
independent choices at two different endpoints.
\end{proof}

\subsection{Entrywise products and difference matrices}

\begin{lemma}[Diagonal prediction and Hadamard reduction]
\label{lem:path-hadamard}
Assertion \textup{(iv)} of \Cref{thm:path-calculus} holds.
\end{lemma}

\begin{proof}
First prove the entry predictions, and then strengthen the resulting
Hadamard estimate to an operator bound.  For independent GOE and
deterministic binned diagonals, each entry of a fixed word is
Frobenius Lipschitz on a spectral cutoff, with a dimension-free
constant.  Gaussian concentration and a union bound over $n^2$
entries therefore give uniform convergence to the entry expectations
at every fixed positive accuracy.
Conjugation by diagonal sign matrices makes all off-diagonal
expectations zero.  A fixed-degree Wick expansion of the diagonal
entry leaves precisely its complete noncrossing contraction with
the outer site type fixed.  Every other pairing has at least one
additional index identification and costs $O(1/n)$, uniformly in the
outer type, including a rare type.  This proves
\Cref{eq:path-entry-prediction}.  The same expansion for the normalized
trace averages the remaining outer diagonal and gives the vacuum
trace rule.

For the off-diagonal Hadamard product, its expectation depends only on
the two endpoint types.  A Wick pairing connecting the two distinct
fixed endpoints has at least one fewer free index than paired edges,
so each mean entry is $O(1/n)$.  If $E_a$ is a type projection, this
mean matrix is, up to an $O(1/n)$ diagonal correction, a sum of dyads
\[
 \sum_{a,b}c_{ab}\,(E_a\1/\sqrt n)(E_b\1/\sqrt n)^t,
 \qquad |c_{ab}|\le C.
\]
It is visible.  We must still prove that the centered product is
small in operator norm; entrywise smallness alone would not prove it.

We obtain the operator bound by treating diffuse and large
coordinates separately.  Start with unit vectors $u,v$ satisfying
$\max(\norm u_\infty,\norm v_\infty)\le\zeta$ and consider
$u^t(Q\circ R)v$.  Perturbing $G$ by $\dot G$ gives terms of
the form
\[
 \ip{D_uR D_v}{\dot Q}_{\HS}
       +\ip{D_uQ D_v}{\dot R}_{\HS}.
\]
The word differentials are Hilbert--Schmidt bounded by
$C\norm{\dot G}_{\HS}$.  Moreover,
\[
 \norm{D_uR D_v}_{\HS}
 \le\norm u_\infty\norm R\norm v\le C\zeta.
\]
Hence this bilinear statistic has Gaussian concentration rate
$\exp(-cn a^2/\zeta^2)$.
Nets of the two unit balls at a fixed resolution have size
$\exp(C_an)$.  Choose $\zeta$ sufficiently small after the required
accuracy $a$.  A union bound controls all diffuse unit tests.

To cover all unit vectors, separate the coordinates of magnitude
greater than $\zeta$.  Each vector has at most $\zeta^{-2}$ such
coordinates.  After removing the diagonal,
\[
 \sum_{j\ne i}|Q_{ij}R_{ij}|^2
 \le\max_{j\ne i}|Q_{ij}|^2\sum_j|R_{ij}|^2=o(1),
\]
and the same holds for columns.  Thus the portions involving the
finitely many large coordinates have norm $o(1)$ uniformly in their
locations.  The finite-net approximation is legitimate because the
Schur bound gives a uniform a priori operator norm.  Combining sparse
and diffuse parts proves that the centered off-diagonal Hadamard
product is $o_{\op}(1)$.

Visible dyads do not spoil these claims.  A bounded path vector $b$
with uniform spatial square tails has
$\max_i|b_i|\to0$, since
\[
 \max_i|b_i|^2
 \le M^2/n+\sum_i|b_i|^2\1_{\{\sqrt n|b_i|>M\}},
\]
with dimension first and then $M\to\infty$.  Thus a visible dyad
has negligible maximum entry.  Its Hadamard product with a bounded
matrix is $D_bQD_c$ and has norm at most
$\norm b_\infty\norm Q\norm c_\infty=o(1)$.
Finite-rank approximation extends the argument to the visible closure.

Finally use \Cref{lem:path-compressed,lem:path-sparse} to transfer
the uniform Gaussian comparisons to actual path labels.  Polynomial
norm differences tolerate the small Frobenius comparison error.
Uniform entry predictions retain the actual outer diagonal on rare
sites.  The dyads of binned type indicators are limits of bounded
path-recipe dyads by tracial approximation.  This proves the stated
Hadamard reduction and diagonal formula on admissible trees.
\end{proof}

\begin{lemma}[Smooth difference and column-flip closure]
\label{lem:path-difference-closure}
Assertion \textup{(vi)} of \Cref{thm:path-calculus} holds.  In particular,
for every fixed smooth site recipe $z$, the matrix
$E(z)_{ij}=d_jz_i$ has bounded operator norm and vanishing off-diagonal
maximum, and belongs to the path-matrix closure.  On scaled vector
recipes its operator bound has the additional factor $n^{-1/2}$.
\end{lemma}

\begin{proof}
We prove closure by induction over the recipe operations.  The spin
input satisfies $E(x)=I$, and a query by $J$ left-multiplies its
difference matrix by $J$.  For the nonlinear step, assume closure
for $z$ and apply a bounded smooth scalar function $\phi$.
Off the diagonal, use $z_i(x^j)=z_i(x)-2x_jd_jz_i$ in Taylor's
formula to obtain
\begin{equation}
 d_j\phi(z_i)
 =\phi'(z_i)E(z)_{ij}
       -x_j\phi''(z_i)E(z)_{ij}^2+R_{ij},
 \qquad |R_{ij}|\le C|E(z)_{ij}|^3.
 \label{eq:path-difference-taylor}
\end{equation}
The remainder's maximum row and column absolute sums tend to zero:
\[
 \sum_j|R_{ij}|
 \le C\max_{k\ne l}|E(z)_{kl}|\sum_j|E(z)_{ij}|^2=o(1),
\]
and similarly for columns.  Keep $d_i\phi(z_i)$ exactly on the
diagonal.  The first two terms in
\Cref{eq:path-difference-taylor} are diagonal multiplications of
$E(z)$ and its Schur square; they belong to the asserted closure by
\Cref{lem:path-hadamard}.  Their operator norms are bounded, and
their off-diagonal maxima vanish.  The exact diagonal is a bounded
site label, obtained from $z_i$ and its already constructed own-spin
difference.  Its smooth limiting versions follow by the same recipe
recursion.  This completes induction over every finite calculation.
Finite sums and products use the exact product rule and the bounded
Schur estimate; no unbounded derivative coefficient appears.

For column flips define $\mathsf F(Q)_{ij}=Q_{ij}(x^j)$.
Starting at the right end of a word, one has the exact recursions
\begin{align*}
 \mathsf F(JQ)&=J\mathsf F(Q),\\
 \mathsf F(D_fQ)
 &=D_f\mathsf F(Q)-2\{E(f)\circ\mathsf F(Q)\}D_x.
\end{align*}
The first uses that $J$ is independent of $x$; the second is the
identity $f_i(x^j)=f_i(x)-2x_jd_jf_i(x)$ with its derivative index
still tied to column $j$.  These recursions express every column-flip
matrix through bounded path words and Schur insertions.  They preserve
the closure just proved.  The conditional derivative of a fixed smooth multitime calculation
requires the full mixed-slot recursion of
\Cref{lem:path-mixed-tensors,lem:path-conditional-derivative}.
After its uniform tails are removed, the equality-network reduction
in \Cref{prop:path-opened-derivative} expresses every opened term
through the matrix closure just described.  Thus the dynamic assertion
uses that additional argument; it is not inferred from first partial
Jacobian bounds alone.
\end{proof}

\subsection{Opened mixed derivatives as directed networks}

To handle conditional derivatives, we still need to express the mixed
tensors as path matrices.  Their existing norm bound does not supply
that expression: a Taylor remainder may have a bounded coefficient
depending on several derivative indices.  We first approximate that
coefficient with control in the flattened operator norm.  The
resulting tensors can then be drawn as finite equality networks,
which we reduce after the observation slots have been copied.

\begin{lemma}[Polynomial approximation in the increments]
\label{lem:path-increment-approximation}
For each fixed smooth calculation of
\Cref{lem:path-mixed-tensors}, every nonempty mixed tensor $T_IF$
is uniformly approximable in the norm of
\Cref{eq:path-mixed-norm} by finite tensor polynomials.  Their factors
are matrix-query edges and equality tensors, their vertex labels are
bounded smooth functions of the calculation values, and every
derivative slot occurs in every nonzero term.  Approximation constants
are dimension independent.  No bounded range of the undifferentiated
linear-query values is required.
\end{lemma}

\begin{proof}
Only the increments need polynomial approximation; the base query
values may remain unbounded.  In the corner expansion from
\Cref{lem:path-mixed-tensors}, the formal increments satisfy
$|t_{b,S}|\le 2^{|S|}C_S(z^b)$ and hence range in a fixed compact
box.  Each remainder has the form
\begin{equation}
 \left(\prod_{\alpha\in A}t_\alpha\right)t_\gamma
                 R_{A,\gamma}(z,t),\qquad\gamma\in A.
 \label{eq:path-increment-remainder}
\end{equation}
Here $z$ consists of the undifferentiated calculation values at the
output site.  The integral formula for $R$ and the global scalar
derivative bounds give a uniform bound and a uniform modulus of
continuity in $t$, independently of $z\in\R^q$.

Apply multivariate Bernstein approximation on the increment box,
treating $z$ as a parameter.  Given $\eta>0$, this gives a fixed
finite polynomial
\[
 P_\eta(z,t)=\sum_\nu b_\nu(z)t^\nu,
 \qquad \sup_{z,t}|R_{A,\gamma}(z,t)-P_\eta(z,t)|\le\eta.
\]
In the Bernstein basis the coefficients are values of $R$ at fixed
increment grid points.  They are therefore bounded smooth functions
of $z$ with bounded derivatives of each required order.  Conversion
to the fixed monomial basis preserves these properties.  This
approximates only the increments, not the possibly unbounded $z$.

The error still contains the repeated factor $t_\gamma$ in
\Cref{eq:path-increment-remainder}.  The repeated-square estimate
in \Cref{lem:path-tensor-product} therefore bounds its flattened
operator norm by $C\eta$.  Every retained monomial still covers
the original set $I$ of derivative slots.  Carry out this operation
at every nonlinear gate, and substitute the already constructed
approximations of its input tensors.  The product bound in
\Cref{eq:path-tensor-product} makes each finite substitution
continuous in those operator norms.  Finite induction proves the
asserted approximation.

We record the spatial structure of the resulting terms for the
conditional derivative calculation.  Unroll one term into a tree
rooted at its output site.  A linear gate gives a matrix edge.  At
a tensor monomial, take copies of the nonempty child tensors and
join their output roots at an equality vertex labelled by $b_\nu$.  Each spin input
supplies an identity edge ending in a derivative leaf; retain this
edge even if it could immediately be contracted.  A constant input
has no nonempty derivative term.  Factors
$x^{(a)}_{j_a}$ from the corner increments become bounded labels
at the corresponding leaves.  The covering property ensures that
every required derivative slot has a leaf.  Several occurrences of
one slot remain several leaves with the same slot name, to be copied
by equality tensors when the temporal indices are attached.  Thus no
coefficient depending on two unrelated vertex indices remains.
\end{proof}

\begin{lemma}[Reduction of finite directed equality networks]
\label{lem:path-directed-network}
Let $\mathcal A$ be a class of uniformly bounded finite matrix
expressions, closed under products, transposes, Hadamard products,
and multiplication by its bounded diagonal labels.  Suppose every
such expression $Q$, including expressions newly formed by these
operations, has
\[
 \max_{i\ne j}|Q_{ij}|=o(1),\qquad
 \max_i|Q_{ii}-\delta_Q(i)|=o(1),
\]
where $\delta_Q$ is an admitted bounded diagonal label.  Consider
a fixed finite equality network with matrix edges from $\mathcal A$
and admitted bounded vertex labels.  If it has an acyclic orientation with
unique source and sink at its external input and output, then its
matrix contraction is, up to $o_{\op}(1)$, in $\mathcal A$.
It in particular has vanishing off-diagonal maximum.
\end{lemma}

\begin{proof}
We first bound a network contraction, then show how to reduce the
network while controlling each error.  Put $H=\R^n$.  For
$p,q\ge1$, the labelled equality gate
$E^d_{p,q}:H^{\otimes p}\to H^{\otimes q}$ sends
$e_a^{\otimes p}$ to $d(a)e_a^{\otimes q}$ and annihilates
nonconstant input tuples.  Its norm is $\norm d_\infty$.
Count the external input as an incoming wire at the source and the
external output as an outgoing wire at the sink.  Each vertex now
has both an incoming and an outgoing wire.  Executing vertices in
topological order, with tensor products on the other live wires,
gives the deterministic bound
\begin{equation}
 \norm{Z_G}_{\op}
 \le\prod_{e\in E(G)}\norm{A_e}_{\op}
                    \prod_{v\in V(G)}\norm{d_v}_\infty.
 \label{eq:path-network-bound}
\end{equation}
In particular, there is no free summation map with norm $\sqrt n$.

Call an internal edge $e:u\to v$ safe if deleting it preserves
the unique source and sink and makes $u,v$ incomparable in the
remaining directed order.  Both endpoints then retain positive
incoming and outgoing degrees.  Process their predecessors first
and process $u,v$ jointly.  If their remaining degrees are
$p_u,q_u,p_v,q_v\ge1$, the joint gate with coefficient $R$ is
\[
 e_a^{\otimes p_u}\otimes e_b^{\otimes p_v}
 \longmapsto d_u(a)d_v(b)R_{ab}\,
       e_a^{\otimes q_u}\otimes e_b^{\otimes q_v},
\]
and is zero on all other input tuples.  Both the displayed input
family and output family are orthonormal.  Its norm is at most
$\norm{d_u}_\infty\norm{d_v}_\infty\max_{a,b}|R_{ab}|$.
Taking $R$ to be the edge matrix restores the original network;
taking its diagonal identifies $a=b$ and contracts $u,v$.  Hence
\begin{equation}
 \norm{Z_G-Z_{G/e}}_{\op}
 \le\max_{a\ne b}|(A_e)_{ab}|
       \prod_{f\ne e}\norm{A_f}_{\op}
       \prod_w\norm{d_w}_\infty=o(1),
 \label{eq:path-network-pinch}
\end{equation}
with new vertex label $d_ud_v\diag(A_e)$.
Its diagonal can be replaced by $d_ud_v\delta_{A_e}$ using
\Cref{eq:path-network-bound}.  Contraction cannot create a directed
cycle: such a cycle would yield a path from $u$ to $v$ or from $v$
to $u$ after deletion.  It preserves the unique source and sink.

The norm estimate justifies pinching a safe edge.  To ensure that
such an edge is available, first consolidate parallel edges by
Hadamard products and suppress every nonterminal vertex of indegree
and outdegree one, using matrix multiplication with its diagonal
label.  Continue until neither reduction applies.  If any internal
vertex remains, the earliest one in a topological order has just
one incoming edge, from the source.  It must consequently have
at least two outgoing edges, or it would have been suppressed.
Choose the latest internal vertex $u$ having at least two outgoing
edges, and let $v$ be its earliest immediate successor.  The
successors are distinct, so $v$ is internal.  It has outdegree one
by the choice of $u$, and indegree at least two because suppression
has ended.  There is no alternate path $u\to v$: its first edge
would go to another successor of $u$, later than $v$, which cannot
subsequently reach $v$.  Deleting $u\to v$ therefore makes them
incomparable and leaves positive incoming and outgoing degrees at
both endpoints.  This is a safe edge.

Pinch that edge and repeat the ordinary reductions.  A pinching or
series reduction decreases the number of vertices, and parallel
consolidation decreases the number of edges.  The procedure therefore
ends after finitely many steps with one edge from source to sink.
All newly formed edges belong to $\mathcal A$, so their own entry
predictions justify every subsequent pinching.  Summing the finitely
many errors proves the claim.  If source and sink have already been
identified, acyclicity leaves only a diagonal equality gate; this
case is immediate.  Finally
$\max_{ij}|E_{ij}|\le\norm E_{\op}$ transfers the off-diagonal
assertion through the accumulated errors.
\end{proof}

Figure~\ref{fig:network-pinching} shows the first crossed example:
pinching one internal edge turns a contraction with two internal
indices into a product of two Hadamard factors and a diagonal label.
The proof above applies the same operation after each finite
series and parallel reduction.

\begin{figure}[htbp]
\centering
\begin{tikzpicture}[
  vertex/.style={circle,draw=black!65,fill=white,inner sep=2pt},
  edge/.style={-{Stealth[length=1.8mm]},draw=black!65},
  lab/.style={font=\small,fill=white,inner sep=1pt}]
 \node[vertex] (j) at (0,0) {$j$};
 \node[vertex] (g) at (1.65,1) {$g$};
 \node[vertex] (k) at (1.65,-1) {$k$};
 \node[vertex] (i) at (3.3,0) {$i$};
 \draw[edge] (j) -- node[lab,above left] {$C$} (g);
 \draw[edge] (j) -- node[lab,below left] {$Q^{\mathsf T}$} (k);
 \draw[edge] (g) -- node[lab,above right] {$B$} (i);
 \draw[edge] (k) -- node[lab,below right] {$E$} (i);
 \draw[edge,line width=1pt] (g) -- node[lab,right] {$T^{\mathsf T}$} (k);
 \node at (4.3,0) {$\longrightarrow$};
 \node[vertex] (jj) at (5.3,0) {$j$};
 \node[vertex] (v) at (7.6,0) {$a$};
 \node[vertex] (ii) at (9.9,0) {$i$};
 \draw[edge] (jj) -- node[lab,above] {$C\circ Q^{\mathsf T}$} (v);
 \draw[edge] (v) -- node[lab,above] {$B\circ E$} (ii);
 \node[font=\small,below=5mm of v] {$T_{aa}$};
 \node[font=\small] at (1.65,-1.65) {two internal indices};
 \node[font=\small] at (7.6,-1.65) {one internal index};
\end{tikzpicture}
\caption{Pinching the bridge in a crossed contraction.
The left network represents
$\sum_{g,k}B_{ig}E_{ik}T_{gk}C_{gj}Q_{jk}$.
Removing the off-diagonal entries of $T$ costs at most a fixed
operator-norm constant times $\max_{g\ne k}|T_{gk}|$.
Its diagonal identifies $g=k=a$, leaving
$(B\circ E)D_{\operatorname{diag}T}(C\circ Q^{\mathsf T})$.
The diagonal label is retained.}
\label{fig:network-pinching}
\end{figure}


\begin{proposition}[Opened conditional derivatives]
\label{prop:path-opened-derivative}
Fix a smooth multitime calculation as in
\Cref{lem:path-mixed-tensors}, a finite causal tree, common
Duhamel and likelihood truncations from
\Cref{lem:path-conditional-derivative}, and the tensor
approximations of \Cref{lem:path-increment-approximation}.
Require every linear-query edge of the calculation and every opened
$K$ edge to belong to the admitted path-matrix algebra, with its
entry predictions, and require all bounded coefficient labels to be
admitted path labels.  In particular, the calculation values are
admitted path recipes; the likelihood calculations above use $J$
and $J\circ J$ as their query matrices.  The broader class of
arbitrary bounded-operator query maps in
\Cref{lem:path-mixed-tensors} is not asserted to have this closure.
After ordinary conditional expectations are opened on their specified
causal branches, every remaining derivative contraction is a bounded
path matrix up to successive operator-norm errors.  It has vanishing
off-diagonal maximum and the Hadamard and diagonal predictions of
\Cref{lem:path-hadamard}.  The network errors have the typical,
weighted and integrated-parameter meaning of \Cref{def:path-tree}.
The earlier tensor and conditional truncation errors are uniform in
the joining state and stored parameters.
\end{proposition}

\begin{proof}
The first step removes errors that are uniform in the joining state.
Use \Cref{eq:path-copy-bound,eq:path-amplified-tail} for the tensor
approximations and the common $K$ tails.  At fixed accuracy, only
finitely many terms and fixed finite-measure Duhamel integrals remain.
Every $K$ insertion history is a bounded finite path-matrix word,
with a fixed cap on each likelihood label.  Keep ordinary $S$
expectations exact for now.  On opening a later inner-product or
replica test, they become ordinary causal branches and introduce no
new spatial summation index.

We now identify the two trees whose gluing represents an opened term.
The spatial recipe tree is the one constructed in
\Cref{lem:path-increment-approximation}.  A second, temporal tree
starts at the joining derivative coordinate: construct it by reading
the finite conditioning recursion backwards.  In
\Cref{eq:path-copy-tensors}, the memory term keeps its current
slot $p$.  The active term replaces $p$ by the newly introduced
slot $m$, through the $K$ path-word edge.  The mixed term forks
into $p$ and $m$, puts the $K$ edge on the latter branch, and
puts the bounded sign $-2x_p$ at the fork.  The new slot $m$ was
absent from the current derivative subset.  It therefore cannot
reconnect to an existing temporal leaf.  Repeating this construction
gives a temporal tree with one terminal for each original derivative
slot used in that term.  Finite causal-tree conditioning has the same
construction, with copies attached at the common joining vertex.

If a slot occurs several times in the recipe polynomial, append an
equality fork to its temporal terminal, one outgoing branch for each
occurrence.  Glue the resulting temporal leaves to the recipe leaves
in pairs, retaining the input-spin identity edges until this gluing
is complete.  Every leaf is paired.  The covering property ensures
that there is at least one leaf, and that both trees are connected.
Orient the temporal tree away from the joining derivative index and
the recipe tree toward its output site.  A directed path can cross
from the temporal tree into the recipe tree only at a paired leaf;
it cannot cross back.  Thus the glued network is acyclic, with its
unique source the joining derivative coordinate and its unique sink
the output site.  Every vertex is on a path between them.  The only
free indices are these two external ones; all internal indices are
ordinary equality contractions with no normalization factor.

The glued graph now has exactly the orientation required by
\Cref{lem:path-directed-network}.  Its matrix algebra, including
bounded visible dyads, is closed under ordinary products and
transposes.  Lemma~\ref{lem:path-hadamard} supplies the Hadamard
reduction and entry predictions for each newly formed fixed
expression.  Applying the network lemma therefore reduces the
opened term to the same path-matrix closure in operator norm.  The bounded scalar coefficients,
spin signs and diagonal predictions are path labels.  Integrate its
pointwise network bound over the fixed Duhamel parameter sets and
use the integrated entry predictions.  No worst auxiliary time or
starting configuration is selected.  The same argument applies to
stationary reversed causal branches.  Actual and predictor terms
keep the same subdivisions and approximations before they are
recombined and centered.

This proves closure for the opened histories used in the subsequent
conditional tests.  It does not assert a uniform-operator-norm
approximation of an arbitrary conditional matrix average by finitely
many sampled histories.  Conditional expectations remain exact
until those branches are opened.  For a still-averaged off-diagonal
maximum one may instead use conditional Jensen:
\[
 \max_{i\ne j}|\E[Q_{ij}\mid X]|
 \le\E[\max_{i\ne j}|Q_{ij}|\mid X].
\]
The uniform errors were removed before this typical network step.
In particular, no spatial square-mean error has been substituted for
an operator error.
\end{proof}

\subsection{Two-root self-averaging and completion of the calculus}

\begin{lemma}[Site histories from independent equilibrium roots]
\label{lem:path-two-root}
For two independent equilibrium samples in the same disorder, and
independent heat-bath site randomness on their trees, the limiting
joint site law of spin, field, and bounded likelihood histories is
the product of their one-root laws.  Consequently stationary site
means of any such bounded history label self-average in spatial
square mean.
\end{lemma}

\begin{proof}
We first separate the two planted starting directions, then propagate
that separation through the finite simulations.  Plant the samples
$\sigma$ and $\sigma'$.  Their overlap $r=n^{-1}\sigma^t\sigma'$
has the independent-uniform-spin law tilted by
$\exp(\beta^2nr^2/2)$, up to a spin-independent constant.  The
Rademacher rate function satisfies $I(r)\ge r^2/2$.
Since $\beta<1$, for every fixed $\varepsilon>0$ the probability
$|r|>\varepsilon$ is at most $e^{-c_\varepsilon n}$.
Conditionally on the spins the Gaussian disorder has the two rank
shifts
$\beta^2(\sigma\sigma^t+\sigma'\sigma'^t)/n$.
The two initial site spins are therefore asymptotically independent,
and the two planted directions are asymptotically orthogonal.

Perform the finite Gaussian recursion from
\Cref{lem:path-regression} for both causal simulations.  Each query
needed for spin and field histories is odd under global spin inversion
of its own root; a score query for likelihoods is also odd.  Fresh
refresh uniforms can be reflected under this inversion without changing
their law.  Deterministic cross-overlap limits of an odd query from
each different tree change sign under either inversion, and hence
vanish.  The Gaussian regression innovations for the two trees have
zero cross covariance and are independent.  By induction their
site recursions are independent, because their initial spins and site
seeds are independent as well.  Likelihood labels are bounded
functions of the resulting histories; their computation needs no
additional noncentered Gaussian query.  Thus their joint limits also
factor.  This induction does not apply to an arbitrary shared
noncentered query such as $J\1$, explaining the restriction in the
statement.

Transfer this factorization to the actual stationary paths using the
two-replica planting comparison and then
\Cref{prop:path-simulation,lem:path-likelihood-approximation}.  For
a bounded label $b$, represent the spatial mean of
$\E[b_i\mid J]^2$ by the empirical product on the two independent
roots.  The factorization makes its limit $(\tau b)^2$, proving
the claimed self-averaging of the site means.
\end{proof}

\begin{proof}[Completion of the proof of \Cref{thm:path-calculus}]
The required empirical space and endpoint predictions are supplied
by \Cref{lem:path-prediction}, together with orthogonality of hidden
vectors.  Lemma~\ref{lem:path-likelihood-approximation} puts the
likelihood labels in this construction with their stated tracial
convention.  The norm comparison then follows from
\Cref{lem:path-fock-norm,lem:path-sparse,lem:path-compressed}: the
representation is on one tree, retains the rare diagonal base, and
uses the Gaussian rate before paying the entropy and density costs.
The entry and Hadamard assertions come from
\Cref{lem:path-hadamard}; the difference and column-flip closure is
\Cref{lem:path-difference-closure}, with the opened conditional terms
just reduced to that closure.  Finally,
\Cref{lem:path-two-root} proves the restricted two-root independence
assertion.

Every construction used the fixed countable dense recipe family and
successive fixed approximations.  Diagonal extraction therefore gives
common high-probability disorder events for any finite initial list
of tests.  Lemma~\ref{lem:path-tree-transfer} permits polynomial
weights and actual conditional copies.  All superpolynomial path
discards were made separately on the original forward marginals.
These observations give the typicality convention and the order of
limits in the theorem.
\end{proof}
